% Source: https://artofproblemsolving.com/community/c3h1777627_solve_for_abc
% Stable identity: aops-topic-1777627-q01
\textbf{A factorisation.} Multiply the equation $abc-5a-3b=4c-19$ by $c$ and add $15$:
\[
(ac-3)(bc-5)=abc^2-5ac-3bc+15=c(4c-19)+15=(c-1)(4c-15).
\]
\textbf{$c=1$.} Then $(a-3)(b-5)=0$: every triple $(3,b,1)$ and every triple $(a,5,1)$ works (check: $15+3b+4-3b=19$ and $5a+15+4-5a=19$).

\textbf{$c=2$.} $(2a-3)(2b-5)=-7$ with both factors odd and $2a-3\ge-1$, $2b-5\ge-3$: the factorisations $(-1)(7)$ and $(7)(-1)$ give $(a,b)=(1,6)$ and $(5,2)$; the others give $b$ or $a$ not positive.

\textbf{$c=3$.} $3(a-1)(3b-5)=-6$, i.e.\ $(a-1)(3b-5)=-2$. Here $a-1\ge0$, and $a-1=0$ is impossible, so $a-1\in\{1,2\}$: $a-1=1$ gives $3b-5=-2$, so $(a,b)=(2,1)$; $a-1=2$ gives $3b-5=-1$, which has no integer solution.

\textbf{$c\ge4$.} Now the right side $(c-1)(4c-15)$ is positive. If $a\ge2$ and $b\ge2$, both factors on the left are positive and
\[
(ac-3)(bc-5)\ge(2c-3)(2c-5)=4c^2-16c+15>4c^2-19c+15=(c-1)(4c-15),
\]
a contradiction. So $a=1$ or $b=1$.
\begin{itemize}
\item $a=1$: $(c-3)(bc-5)=(c-1)(4c-15)=(c-3)(4c-7)-6$, so $c-3\mid6$ and $c\in\{4,5,6,9\}$. Then $bc-5=(4c-7)-\dfrac{6}{c-3}$ equals $3,10,15,28$, giving $b=2$ for $c=4$, $b=3$ for $c=5$, and no integer $b$ for $c=6$ or $c=9$.
\item $b=1$: $(ac-3)(c-5)=(c-1)(4c-15)=(c-5)(4c+1)+20$, so $c\ne5$ and $c-5\mid20$, giving $c\in\{4,6,7,9,10,15,25\}$. Then $ac-3=4c+1+\dfrac{20}{c-5}$ equals $-3,45,39,42,45,63,102$, giving $a=0$ (not allowed) for $c=4$, $a=8$ for $c=6$, $a=6$ for $c=7$, $a=5$ for $c=9$, and no integer $a$ for $c=10,15,25$.
\end{itemize}
\textbf{Answer.} The solutions are the two infinite families
\[
\boxed{(3,b,1)\ (b\ge1)\quad\text{and}\quad(a,5,1)\ (a\ge1)}
\]
together with the eight triples
\[
\boxed{(1,6,2),\ (5,2,2),\ (2,1,3),\ (1,2,4),\ (1,3,5),\ (8,1,6),\ (6,1,7),\ (5,1,9)}.
\]
